%======================================================================
% 1. INTRODUCTION
%======================================================================
\section{Introduction}
\label{sec:intro}

The external geometry of a ballistic entry capsule settles most of its entry
performance before a single trajectory is flown. Nose radius, shoulder radius
and afterbody cone half-angle set the force coefficients and the trim and
stability characteristics, and through them the integrated heat load, the
attainable ground range and the enclosed payload volume. Shape optimization
therefore resolves a three-way trade between thermal load, mobility and
volumetric efficiency, under limits on peak heating, deceleration and static
stability. In most of the literature this trade is resolved deterministically:
one nominal entry state and one aerodynamic model are assumed, and dispersion
is absorbed downstream by margins in the thermal-protection and guidance
design, without feedback into the shape itself.

Deterministic capsule shape optimization is well explored. Dirkx and Mooij
\cite{dirkx2014,dirkx2017entry} optimize capsule shapes with
local-inclination aerothermal models and evolutionary algorithms; the
present simulation chain descends from that work. Related aerodynamic
capsule designs and parameterized-shape optimizations are reported by Chen
et al.\ \cite{chen2015}, and the surrogate-modelling toolbox behind such
workflows \cite{forrester2008surrogate} rests on the
design-and-analysis-of-computer-experiments literature
\cite{sacks1989,jones1998}.
The uncertainty-aware literature is thinner. Ridolfi et al.\
\cite{ridolfi2012robust} made a multi-disciplinary capsule optimization
robust to uncertainty in design parameters and mission conditions, and
Ghaedamini Harouni and Hashemi Mehneh \cite{hashemi2019capsule} solved a
robust multi-objective capsule shape problem --- maximizing volumetric
efficiency among the objectives --- with adaptive Monte-Carlo inner loops,
reporting constraint satisfaction at a 99.8\% confidence level. Both
establish that robustness can be enforced on capsule shape design; neither
compares the robust result against a deterministic baseline solved on the
same physics, which is what isolates the effect of the formulation itself.
The general theory of robust optimization \cite{beyer2007robust} and the
reliability-based tradition surveyed by Valdebenito and Schu\"eller
\cite{valdebenito2010} flag two recurring obstacles that are visible in
these applications: the cost of the inner uncertainty propagation, and the
weak evidence base of the uncertainty model itself.

For capsule shape design specifically, three gaps remain. First, the
uncertainty bands entering robust or reliability-based formulations are
usually assumed: entry-state dispersions are inherited from targeting
budgets, and aerodynamic model-form error is assigned a round percentage with
little justification. Second, deterministic and robust formulations are
rarely solved on the same validated physics, so any reported ``benefit of
robustness'' mixes the formulation change with model and solver differences.
Third, the cost of ignoring uncertainty is seldom quantified, and even less
often verified: what a deterministic front hides in constraint-exceedance
probability and objective dispersion is typically not measured against the
true model.

This paper closes the three gaps with a controlled comparison on a single
simulation chain. The uncertainty model is evidence-based: entry-state
dispersion (flight-path angle, entry speed) is combined with a multiplicative
aerodynamic model-form uncertainty whose 5\% standard deviation is sized from
the residuals of a validation campaign against Apollo wind-tunnel and flight
data, not chosen for convenience. Deterministic and robust formulations of
the same three-objective problem are then solved on the same Gaussian-process
surrogates, evaluated by a common-random-number moment engine so that the two
fronts are directly comparable, and the deterministic front is additionally
audited under the full uncertainty model. The comparison is quantified rather
than asserted: at matched expected performance, robust designs cut the
standard deviation of integrated heat load in 98.1\% of 5{,}371 design pairs
(pair-level median 27.5\%; 35.0\% against the robust front's
$\sigma$-envelope at matched performance), while the deterministic front
conceals
shoulder-heating exceedance probabilities up to 10\%, confirmed at 9.7\% by
Monte-Carlo simulation of the full model. An irreducible dispersion floor of
\SI{16}{\mega\joule\per\square\metre} (coefficient of variation about 20\%)
bounds both fronts and is set by the flight-path-angle dispersion, not by
geometry. Finally, selected designs are verified against the trajectory
simulator itself, 1{,}500 runs in total; the same exercise exposes a
systematic bias of the mixture surrogate in predicted range, which we
disclose and bound.

Section~\ref{sec:model} describes the vehicle model and its validation.
Sections~\ref{sec:surrogate} and \ref{sec:formulation} present the
uncertainty model, the surrogate construction and its accuracy gate, and the
two design formulations. Section~\ref{sec:results} reports the front
comparison, Section~\ref{sec:verification} the Monte-Carlo verification, and
Sections~\ref{sec:discussion} and \ref{sec:conclusions} the discussion and
conclusions.
